\PBchapter{Vida, viagens e confusões}
Devem existir muitos textos filosóficos que definem o que é viver. 

Esses textos devem refletir sobre a existência, sobre a felicidade, sobre relacionamentos, sobre plenitude. 

Nesse momento, eu estou definindo vida sendo algo que ocorre entre os vazios. 

Eu preenchi minha jornada com relatos do que vivi, e vejam que eu conto anos chaves, tais como 1989, 1994, 2002, 2012 e 2019 por exemplo. 

O que tem nesses intervalos? Porque não existem relatos sobre eles? Será que é somente para o texto não ficar extenso e maçante?

Aqui podemos fazer um paralelo com as redes sociais atuais. 

Nelas só colocamos aquilo que foi bom: o chopp no bar, o exercício na academia, a festa, o acampamento, a viagem, o início de um relacionamento. 

Pois bem, no meu caso, eu fiz ao contrário. Relatei para vocês algumas adversidades, porque nesses intervalos é que minha vida acontecia. 

Em outubro de 1996, descubro que posso viajar. 

Sim, parece engraçado. 

Mas para algumas pessoas isso não era algo acessível. Acredito que ainda não seja.

O motivo foi descobrir o bom e velho rock’n’roll. 

Ficamos sabendo que a banca AC/DC viria ao Brasil. 

Conseguimos quatro ingressos, então partimos para Curitiba. 

Precisam entender, estamos falando de alguém com 19 anos, e que mal tinha saído por conta própria de Cascavel. 

Foi muito marcante, e achei mesmo que a partir daquele momento ia aprender guitarra e virar um músico. 

Mas como já dei spoiler, é lógico que isso não aconteceu. 

Essa viagem teve momentos icônicos e que guardo na memória até hoje. 

O show foi algo surreal.

Em dezembro do mesmo ano, viajei para a praia com mais quatro colegas da empresa. Conheci o mar. 

Viagem dessas, decidida de última hora. 

Sem lugar pra ficar. 

Sabendo somente a cidade. Piçarras-SC.
 
Mais uma viagem com momentos marcantes, com histórias contadas por mim até hoje.

Janeiro de 1998, vestibular, matemática, Unioeste. 

Não sei como descrever como foi passar nesse vestibular. Nessa época era um acontecimento.

Eu estava no trabalho, o rádio, sim, rádio, ligado. 

Então começam a anunciar os aprovados do curso de matemática, e pasmem, em ordem de classificação. 

Eram 40 vagas, e o curso era relativamente concorrido na época. 

Quando anunciaram os 35 primeiros, eu já tinha desistido. 

Porém, eis que o locutor anuncia o 39o. e na minha cabeça, ele dá uma pausa de uns 5 segundos, e então anuncia meu nome. 

Meus colegas de trabalho comemoram. Eu fico em um estado catatônico, sem reação. Acredito que poucos vão ter a noção de como isso foi uma vitória.
	
Mas aí, chega março de 1998, as aulas começam. 

E sim, o calendário estava todo bagunçado de greves anteriores. 

E eu percebi que passar no vestibular, foi o mais fácil. Trágico. Final de Maio, a primeira prova de Cálculo, nota 0,2. 

O cara que era tido por muitos como inteligente, o melhor em matemática no ensino médio. Então pela primeira vez, percebi que eu teria que estudar. 

É, quem diria, que no ensino médio e fundamental ensinam tanta coisa, menos a estudar. 

Mas consigo recuperar e passar em cálculo, mas foi um ano que ninguém sabia de mim, até Dezembro de 1998.
 
Como me esforcei muito para passar em cálculo, acabei reprovando em duas outras disciplinas, que nem eram tão difíceis.

Mas enfim, em Dezembro de 1998, mais uma viagem para a praia. Mesma cidade, Piçarras-SC, agora com 16 pessoas. 

Épico. Histórico. 

Brigas. Rolo. Confusão. 

Aqui muita Vida aconteceu.

Começa 1999, e eu envolvido com Centro Acadêmico de Matemática, e já me metendo com o Diretório Central de Estudantes. 

Depois de anos com chapa única, pela primeira vez, ia ter uma concorrência.  

Chapa Mostra a Tua Cara! As reuniões eram na minha casa.

Mais uma vitória
. 
Esse ano teve muitos acontecimentos. 

Teve Carnaval de salão, não sei se é comum em outras cidades. 

Grito de carnaval, Carnaval normal e Carnaval de inverno. 

Teve excursões para bagunças e festas e para estudos. 

Teve mais um pouco de vida. 
	
Em outubro de 2009, mais um show do AC/DC, na verdade, dois shows. Buenos Aires-ARG.

Momento louco. Conhecer outro país e ainda curtir um show de rock. 

A vida também aconteceu aqui.

Viajei para as capitais, Assunção, Buenos Aires (novamente), Santiago e Montevideo. 

Acampamentos.

Viagem de carro. Sete pessoas. Dois carros. De Cascavel até Natal-RN.

A vida também aconteceu aqui.

Sei que muitos, nesse momento, vão pensar: “mas isso muitos fazem, as pessoas viajam”

As pessoas vão à praia. 

As pessoas vão a shows. 

Gostaria que refletissem sobre o meu relato. 

Pensem novamente da onde eu vim. 

Para fazer algo assim, em meio a tantos sobressaltos, você precisa querer muito. 

Se dedicar. 

Outros vão pensar: 

“Por isso que está ferrado financeiramente”.

“Por isso não consegue encontrar um propósito”

É uma forma de se pensar, e eu aceito. 

Mas quantas pessoas chegam no final de suas vidas, com suas casas, terrenos, apartamentos, carros, e sem uma história sequer para contar? 

Sem ter sentido a brisa do mar. 

Sem ter passado um perrengue em uma viagem.

Sem falar Sorviete de Muriango, tira la película, Picture.

Sem ter a verdadeira sensação de estar vivo.